\section{評価実験}\label{sec5}

本章では，提案手法のタスク遂行能力と各構成要素の働きを，一連の実験と分析によって評価する．基本性能評価では状態変化312件・配置103件を対象に，提示方式・モデル・前提条件検証の有無による成功率と行動列長を調べる．構成要素比較では，代表4タスクの240エピソードを用いて，検証方式，失敗後修復，再考，事例検索の違いを比較する．さらに，固定入力での判定評価とケース分析によって，どの判断や実行過程に失敗が残るかを検討する．

各評価は同じタスク集合を基礎とするが，対象件数，初期化条件，記録の粒度は異なる．基本性能の18条件（各タスク種別9条件）と構成要素比較の結果はそれぞれの分母で示し，両者を合算した成功率は求めない．代表候補の再実行と人手確認は，個別過程を解釈するための分析として位置付ける．

\subsection{評価用データセット}\label{sec5.1}
本研究では提案手法を評価するために，\ref{sec3}で示した要件を満たす新たなデータセットを構築した．具体的には，物体の状態に着目した「状態変化タスク」と，物体間の関係に着目した「配置タスク」の2種類のタスクを構築した．

\subsubsection{状態変化タスク}
評価データに含まれる状態変化の指示は``Turn on all lightswitches''のようなON/OFF操作である．開閉操作は計画器の許可操作に含まれるが，窓や扉の開閉を最終ゴールとする指示は今回の状態変化データに含まれない．対象物の全種類と件数は\ref{tab:dataset_classes}に示す．

各指示は複数の同名物体を対象とし，同じ指示に初期状態の異なるインスタンスを対応させる．すべて目標状態の場合と一部が未達の場合とでは，必要な行動が異なる．具体的な件数・参照列長を\ref{tab:dataset_extent}に，必要な判断を\ref{sec5.cases}に示す．

\subsubsection{配置タスク}
本タスクでは``Put all apples in the fridge''などの指示に対して，複数の物体を環境内の特定の場所に配置することが求められる．対象物には，ロボットが把持可能な物体と，他の物体を受け入れることができる物体（``receptacle''）が含まれる．

複数物体の保持・配置と残り対象の追跡を扱う．冷蔵庫のような開閉可能な配置先には，初期状態がOPENとCLOSEDのインスタンスがあり，格納前の開扉の要否を判断する必要がある．一方，すべての配置先が開閉可能な容器とは限らず，ON関係をゴールとする支持物も含む．

\subsubsection{データセット構築}
配置タスクではVirtualHome内の移動可能な物体を，状態変化タスクでは状態変更可能な物体を対象として選定した．状態変化タスクの指示は，選定した物体のON／OFF状態の変更を扱う．

\begin{figure}[tb]
    \begin{lstlisting}[caption=状態変化タスクのデータセット例, label=state-change-task]
{
    "task": "Turn on all lightswitches",
    "scene": 1,
    "initial_room": "kitchen",
    "initial_states": [
        {"id": 71, "states": ["ON"]},
        {"id": 173, "states": ["OFF"]},
        {"id": 261, "states": ["ON"]},
        {"id": 427, "states": ["OFF"]}
    ],
    "goal_states": [
        {"id": 71, "states": ["ON"]},
        {"id": 173, "states": ["ON"]},
        {"id": 261, "states": ["ON"]},
        {"id": 427, "states": ["ON"]}
    ],
    "action_scripts": [
        "[WALK] <lightswitch> (173)",
        "[SWITCHON] <lightswitch> (173)",
        "[WALK] <lightswitch> (427)",
        "[SWITCHON] <lightswitch> (427)"
    ]
}
    \end{lstlisting}
\end{figure}

\begin{figure}[tb]
    \begin{lstlisting}[caption=配置タスクのデータセット例, label=placement-task]
{
    "task": "Put all plums in the fridge",
    "scene": 4,
    "initial_room": "bedroom",
    "initial_states": [
        {"id": 103, "states": ["CLOSED"]}
    ],
    "goal_states": [
        {
            "from_id": 53,
            "to_id": 103,
            "relation_type": "INSIDE"
        },
        {
            "from_id": 54,
            "to_id": 103,
            "relation_type": "INSIDE"
        }
    ],
    "action_scripts": [
        "[WALK] <plum> (53)",
        "[GRAB] <plum> (53)",
        "[WALK] <fridge> (103)",
        "[OPEN] <fridge> (103)",
        "[PUTIN] <plum> (53) <fridge> (103)",
        "[WALK] <plum> (54)",
        "[GRAB] <plum> (54)",
        "[WALK] <fridge> (103)",
        "[PUTIN] <plum> (54) <fridge> (103)"
    ]
}
    \end{lstlisting}
\end{figure}

本研究では，\ref{sec4.6}で述べたように，シミュレータとしてVHを使用する．したがって，データセットはVHとの互換性を持つように設計した．

VHには7種類の異なるシーン（家屋）があり，それぞれに複数の部屋とそこに設置された様々な物体が含まれている．各タスクは，これらシーン毎の特性に応じて構成している．

本データで扱う主な状態はON，OFF，OPEN，CLOSEDである．状態変化の指示の対象は7種類，配置の対象は12種類，配置先は9種類である．\ref{tab:dataset_classes}に，指示文中の表記と全インスタンス数を列挙する．物体クラス名との照合には\ref{sec4.1}の単数形補完を用いる．配置先はタスクに明示された対象である．

\begin{table*}[tb]
\caption{指示文中の対象・配置先の全種類と件数（事例用とテストの合計）}\label{tab:dataset_classes}
\centering\small
\begin{tabular}{p{.18\textwidth}p{.74\textwidth}}
\hline
区分 & 種類（件数） \\
\hline
状態変化対象：7種類 & candles (20), cellphones (44), computers (24), faucets (56), lightswitches (256), tablelamps (40), tvs (24) \\
配置対象：12種類 & alcohols (3), apples (12), bananas (18), bellpeppers (6), chips (14), crackers (6), creamybuns (8), cupcakes (26), juices (5), milks (13), peaches (21), plums (22) \\
配置先：9種類 & coffeetable (9), cuttingboard (1), desk (3), fridge (48), fryingpan (1), kitchencounter (7), kitchentable (26), microwave (56), oventray (3) \\
\hline
\end{tabular}
\end{table*}

データセットの例を，Listing~\ref{state-change-task}および~\ref{placement-task}に示す．データセットは，タスクの記述，シーン，エージェントの初期位置，物体の初期状態，物体の目標状態，そしてタスクを完遂するために必要なアクションスクリプトで構成される．

Listing~\ref{state-change-task}は状態変化の\texttt{test\_task6}，Listing~\ref{placement-task}は配置の\texttt{test\_task65}に対応する（空のcommentフィールドのみ省略）．後者の参照行動列は9ステップであり，冷蔵庫を閉める操作は含まれない．ゴールはID 53と54がID 103のINSIDEにあることで，冷蔵庫がCLOSEDであることは要求しない．途中で毎回冷蔵庫を閉める必要性も評価していない．

部屋・シーンはVirtualHomeが提供するものをそのまま使用し，シーン自体には変更を加えていない．同名インスタンスでは，使用するシーンとエージェントの初期部屋を選び，物体の初期状態配列等を設定する．初期化は提供されたシーンの物体位置を基礎とし，\texttt{initial\_states}に指定した状態配列を上書きする．対象・配置先別の件数は表\ref{tab:dataset_classes}に，シーン・初期部屋別の件数は公開資料\path{Paper/analysis/dataset_dimension_counts.csv}に示す．これらは観測した組合せの集計である．初期状態とゴールは各タスクのデータに明示し，状態変化では対象物体の目標状態を，配置では対象物体と配置先のON／INSIDE関係をゴールとして設定する．

\begin{figure*}[tb]
\footnotesize
\begin{align}
 \label{formula5.1}
 Success Rate(SR) &= \frac{Sum Of Successes}{Sum Of Tasks}\\
 \label{formula5.2}
 Action Execution Failure Rate(AEFR) &= \frac{Sum Of Action Execution Failures}{Sum Of Tasks}\\
 \label{formula5.3}
 Failure Rate Of Reaching Maximun Attempts(FRRMA) &= \frac{Sum Of Reaching Maximun Attempts}{Sum Of Tasks}\\
 \label{formula5.4}
 Erroneous Terminate Failure Rate(ETFR) &= \frac{Sum Of Erroneous Terminate Failures}{Sum Of Tasks}\\
 \label{formula5.5}
 Failure Rate(FR) &= AEFR + FRRMA + ETFR\\
 \label{formula5.6}
 Average Steps &= \frac{Sum Of Action Steps For All Tasks}{Sum Of Tasks}
\end{align}
\end{figure*}


\subsubsection{データの分割方法}
全618件の対象・指示・シーン・状態上書き・初期部屋の対応は，公開データの各インスタンスと\path{Paper/analysis/task_inventory.csv}に記録する．
データの分割方法について説明をする．状態変化タスクデータセットの事例数は464個，配置タスクデータセットの事例数は154個であり，テスト用と訓練用データセットの事例数が大体2対1の割合になるように，各データセットを分割する．

訓練用データセットは，~\ref{sec4.3}で述べた，プロンプトに与えるタスク例に対するアクションスクリプトの出力事例として用いる．
状態変化タスクのテスト用事例数は312個，訓練用事例数は152個，配置タスクのテスト用事例数は103個，訓練用事例数は51個とした．

同名でも異なる初期条件を持つデータを別インスタンスとして扱う．状態変化ではテスト6名・事例用6名，配置ではテスト38名・事例用19名の指示があり，いずれも両集合間の同名指示の重複は0である．一方，両集合にはともにシーン1～7が含まれ，物体クラスや操作構造も共有する．分割は手作業で行い，各集合の件数を調整するとともに，同名指示が両集合に重複しないようにした．評価にはこの分割を記録した公開JSONを固定して用いる．「訓練用」はモデルのパラメータ学習を意味せず，プロンプトに与える事例の候補集合を指す．

例えばテストの``Turn on all lightswitches''に対し，事例集合には``Turn on all tablelamps''がある．配置では``Put all bananas in the fridge''と``Put all apples in the fridge''のように対象だけが異なる指示が両集合に存在する．これは両集合間で操作構造が類似する例である．検索入力はタスク名の埋め込みで，選択規則は\ref{sec4.3}に示した．

事例検索の特性を調べるため，テスト指示と選択事例のコサイン類似度および対象・操作・配置先の一致率を分析した．テストインスタンスで重み付けした平均コサイン類似度は状態変化0.624，配置0.781だった．配置103件での対象クラス・操作・配置先の一致率はそれぞれ89.3\%，54.4\%，10.7\%である．例えば``Put all plums in the fridge''には``Put all plums on the kitchentable''が類似度0.761で選ばれ，``Turn off all tablelamps''には``Turn on all tablelamps''が0.910で選ばれた．この分析では候補名をソートして検索を実行しており，基本性能評価の各試行で採択された事例との一致は，個別検索ログがないため確認できない．構成要素比較で用いる事例名・選択規則と，意味類似検索の類似度・上位差・列長は\ref{sec5.control_settings}に示す．

\begin{table}[tb]
\caption{テスト集合の規模と参照行動列長}\label{tab:dataset_extent}
\centering\small
\begin{tabular}{lrr}
\hline
 & 状態変化 & 配置 \\
\hline
インスタンス数 & 312 & 103 \\
異なる指示数 & 6 & 38 \\
ゴール要素数の範囲 & 2--5 & 2--6 \\
参照列長（平均） & 3.872 & 9.175 \\
参照列長（中央値） & 4 & 8 \\
参照列長の範囲 & 0--10 & 0--25 \\
参照列長0の件数 & 27 & 2 \\
\hline
\end{tabular}
\end{table}

\ref{tab:dataset_extent}はテスト集合の属性を示す．参照列長0のデータは初期状態で完了しているというデータ側の設定に対応し，終了判定の振る舞いも評価する．参照列は最短解と証明されておらず，ゴール数や列長だけから人間の成功率・難易度を決めることはできない．

\subsection{実験条件と評価方法}\label{sec5.2}
実験の前段階として，データセットのシーンと初期状態を用いてVH環境の初期設定をする．次に提案手法にデータセットのタスク記述を入力して，タスクを実行するための適切な行動計画を生成する．

基本性能評価の提示方式は，\ref{sec4.3}で定義したSingle PromptとMulti Promptsである．前者は要素を連結した一つのuserメッセージ，後者は五つのuserメッセージを用い，いずれも行動生成1ステップにつき1回の呼び出しを行う．前提条件検証あり／なしを各方式で比較する．両方式には終了判定・行動生成の文面とsystemメッセージ末尾の空白にも差がある．検証方式・再考・検索を個別に変える条件は\ref{sec5.additional}で述べる．

アクションステップの生成時に用いる，タスクを遂行するためのアクションスクリプトの生成例の選択方法について説明する．例の選択は，訓練データセットの中から，入力となるタスク名と類似度の高いタスク名を1種類選択する．タスク名に対して複数個のデータが存在する場合には，タスクを遂行するためのアクションスクリプトの生成例が最長のものを選択する．

\ref{sec4.2}の試行上限$M$は，評価対象自身の参照行動列長$L_{\mathrm{ref}}$を用いて$M=2L_{\mathrm{ref}}$とする．検索で選んだ事例の列長ではない．これは参照解の長さを利用する評価用予算である．上限判定は成功実行の直後に行い，$n\geq M$なら，その時点でゴールを満たしていても上限失敗とする．この場合にLLMが追加行動を生成しようとした事実までは観測していない．$L_{\mathrm{ref}}=0$の特例と想定外の終了判定文字列の扱いもAlgorithm~\ref{alg:public_run}に従う．

ベースラインは，Huangらの手法~\cite{huang22}の著者公開実装\footnote{\url{https://github.com/huangwl18/language-planner}}を変更せず，そのまま使用し，本研究の状態変化・配置タスクデータセットで実行した比較条件である．ベースラインの事例選択（Dynamic Example）も公開実装のまま用い，提案手法の事例選択処理への置換や変更は行っていない．両手法とも自然言語指示文であるタスク名をベクトル化して検索する．公開実装では各事例のタスク名とのコサイン類似度から事例を直接選択する．提案手法では，\ref{sec4.3}に記載した通り，異なるタスク名から最も類似する名前を選択した後，同名事例から最長の参照行動列を採択する．掲載値は本研究で得た実験結果である．物体IDの探索はシミュレータの自動探索に委ねる条件とした．

ベースラインのAction Translationは，Planning LMが生成した自由文の行動候補を，Translation LMによる埋め込みのコサイン類似度で許可行動一覧（\texttt{available\_actions.json}）へ対応付ける処理である．公開コードでは，複数の生成候補について，対応先との最大類似度と生成トークンの平均対数確率の重み付き和を計算し，直前と同じ行動には減点する．最高得点の変換後行動を選び，終了条件に該当しなければ履歴に追加して次の生成を行う．この変換は行動表現を許可行動へ対応付けるもので，現在状態における前提条件の充足判定とは異なる．提案手法では，\ref{sec4.3}のプロンプトに許可操作とID付き出力形式を与えて行動を直接生成する．

評価するLLMはGPT-4o miniとGPT-4oの2種類で，ベースラインはGPT-4oのみである．基本性能評価は，2モデル・2提示方式・検証あり／なしの8条件とベースラインの計9条件を各タスク種別に適用する．モデル版と実行条件の確認範囲を以下に示す．

\subsubsection{評価間の条件差と記録範囲}\label{sec5.provenance}
基本性能評価では初期部屋と状態配列を指定し，構成要素比較ではさらに人物位置を固定して参照グラフとの一致を確認する．また，構成要素比較では各判定の入力・応答と実行前後の状態を記録し，途中の失敗実行数やAPI要求数も分析する．

基本性能評価のモデルは\texttt{gpt-4o-mini-2024-07-18}と\texttt{gpt-4o-2024-08-06}を用いた．構成要素比較ではこれらのモデルIDとtemperature 0を要求ログで確認できる．基本性能評価のベースラインは著者公開実装を変更せず，そのまま使用し，本研究のデータセットで実行した．基本性能の提案手法のコード・データ・結果JSONはコミット\texttt{2761158}に含まれ，READMEのみが異なる\texttt{de1fc8d}と照合した．全18条件のSR・3失敗率・平均列長の計90値は，保存結果からの再集計と小数第3位まで一致した．


公開Notebookのモデル指定は\texttt{gpt-4o-mini}という別名であり，Multiの評価ループには1件で停止する\texttt{break}がある．ベースラインのAction Translationには公開実装の処理を用い，その手順は前述の「実験条件と評価方法」に記載した．保存結果の出所と確認範囲は\path{Paper/revision-completion/Baseline-Provenance.md}と\path{Paper/revision-completion/Archive-Resolution.md}で公開する．

\subsection{評価指標}\label{sec5.3}

本研究では，既存研究\cite{ALFRED20,song2023llm}の評価で用いられているタスクの成功率によって提案手法を評価する．
本評価の通常採点では，タスクJSONに指定された物体IDを用いてゴール状態を照合する．

成功率は，式\ref{formula5.1}に示すように，タスクの総数に対する成功したタスクの割合として定義する．公開実装の通常終了では，状態変化は指定IDの状態配列の完全一致，配置は指定した\texttt{from\_id, to\_id, relation\_type}の辺の存在を確認し，すべてのゴールを満たせば成功とする．ゴール以外のノード状態や辺は採点しない．実行失敗と上限到達は通常のゴール採点より優先し，参照列長0・最初のEndでは採点を省略して成功を返す特例がある．

失敗には三つのタイプがある．一つ目はシミュレータにおけるアクション実行による失敗，二つ目は設定された最大試行回数に達したことによる失敗，三つ目はタスク完了と判定したが，タスクのゴール状態をすべて満たせていないことによる失敗である．失敗の原因を分析するために，各失敗タイプの割合を式~\ref{formula5.2}から式~\ref{formula5.4}に定義する．

式\ref{formula5.2}は，シミュレータにおけるアクション実行の失敗に起因する失敗率を示しており，タスクの総数に対するアクションの実行に失敗した回数の割合として定義する．式\ref{formula5.3}は，設定された最大試行回数に達したことによる失敗率を示しており，タスクの総数に対する当該失敗回数の割合として定義する．式\ref{formula5.4}は，タスク完了と判定されたにも関わらず，タスクのゴール状態を完全には満たしていないことによる失敗率を示しており，タスクの総数に対する当該失敗回数の割合として定義する．

また，式~\ref{formula5.6}により``$AverageSteps$''を求める．基本性能評価では結果JSONの\texttt{attempts}を平均し，全18条件で記録行動列長と一致する．提案実装は実行成功した行動のみを記録し，実行失敗した最後の試行や再生成回数，LLM呼び出し回数は含めない．ベースラインについても記録列長との一致は確認したが，記録処理の詳細は不明である．失敗による途中終了でも列長は短くなるため，全件平均の短さだけで効率の優劣を判断しない．参照列長0を含む件数と長さは\ref{tab:dataset_extent}に示した．構成要素比較では，記録行動列長に加えて途中の失敗を含む実行試行数とAPI要求数を区別する．


\subsection{基本性能評価}\label{sec5.4}
\begin{table*}[tb]
    \caption{「状態変化タスク」の評価結果  ※PreCheck:「前提条件の検証」}
    \label{result1}
    \centering
    \begin{tabular}{llccccc}
    \hline
         \multirow{2}{*}{Method} & \multirow{2}{*}{LLM} & \multirow{2}{*}{SR} & \multicolumn{3}{c}{FR} & \multirow{2}{*}{Average Steps}\\
          & & & AEFR & FRRMA & ETFR\\
    \hline
         Baseline~\cite{huang22} & GPT-4o & 0.032 & 0.000 & 0.535 & 0.433 & 10.481\\
    \hline    
         \multirow{2}{*}{Single Prompt} & GPT-4o mini & 0.625 & 0.026 & 0.231 & 0.119 & 4.747\\
         & GPT-4o & 0.788 & 0.016 & 0.173 & 0.022 & 4.381\\
         
         \multirow{2}{*}{Multi Prompts} & GPT-4o mini & 0.750 & 0.045 & 0.183 & 0.022 & 4.917\\
         & GPT-4o & \textbf{0.894} & 0.019 & 0.000 & 0.087 & 3.391\\
    \hline
        \multirow{2}{*}{Single Prompt w/o PreCheck} & GPT-4o mini & 0.615 & 0.003 & 0.253 & 0.128 & 4.872\\
         & GPT-4o & 0.760 & 0.016 & 0.189 & 0.035 & 4.449\\
         
         \multirow{2}{*}{Multi Prompts w/o PreCheck} & GPT-4o mini & 0.696 & 0.016 & 0.253 & 0.035 & 5.080\\
         & GPT-4o & 0.872 & 0.010 & 0.016 & 0.103 & 3.478\\
    \hline
    \end{tabular}
\end{table*}

\begin{table*}[tb]
    \caption{「配置タスク」の評価結果  ※PreCheck:「前提条件の検証」}
    \label{result2}
    \centering
    \begin{tabular}{llccccc}
    \hline
         \multirow{2}{*}{Method} & \multirow{2}{*}{LLM} & \multirow{2}{*}{SR} & \multicolumn{3}{c}{FR} & \multirow{2}{*}{Average Steps}\\
          & & & AEFR & FRRMA & ETFR\\
    \hline
         Baseline~\cite{huang22} & GPT-4o & 0.000 & 0.010 & 0.573 & 0.417 & 15.398\\
    \hline    
         \multirow{2}{*}{Single Prompt} & GPT-4o mini & 0.466 & 0.146 & 0.029 & 0.359 & 7.282\\
         & GPT-4o & 0.738 & 0.117 & 0.029 & 0.117 & 8.078\\
         
         \multirow{2}{*}{Multi Prompts} & GPT-4o mini & 0.408 & 0.117 & 0.029 & 0.447 & 6.447\\
         & GPT-4o & \textbf{0.816} & 0.117 & 0.019 & 0.049 & 8.272\\
    \hline
        \multirow{2}{*}{Single Prompt w/o PreCheck} & GPT-4o mini & 0.476 & 0.117 & 0.019 & 0.388 & 6.592\\
         & GPT-4o & 0.680 & 0.165 & 0.029 & 0.126 & 7.320\\
         
         \multirow{2}{*}{Multi Prompts w/o PreCheck} & GPT-4o mini & 0.437 & 0.126 & 0.010 & 0.427 & 6.136\\
         & GPT-4o & 0.796 & 0.117 & 0.029 & 0.058 & 7.767\\
    \hline
    \end{tabular}
\end{table*}

提案手法とベースライン手法の比較結果について，状態変化タスクデータセットに基づく結果を\ref{result1}に，配置タスクデータセットに基づく結果を\ref{result2}に示す．

\ref{result1}と\ref{result2}の``Baseline''はHuangらの手法~\cite{huang22}の著者公開実装を変更せず，そのまま使用し，本研究のデータセットで実行して得た値であり，引用は手法の出典を示す．``Single Prompt''と``Multi Prompts''は\ref{sec4.3}に述べたメッセージの構成方式を指す．``w/o PreCheck''は，提案手法から前提条件検証とそれに伴う再生成を除いた条件である．

基本性能評価では，提案手法はベースライン条件より高い成功率を示した．条件間では，入力される環境情報・ID情報やベースライン実装の確認範囲が異なる．Multi Promptsは多くの条件でSingle Promptを上回るが，GPT-4o miniの配置タスクでは逆転する．

ベースラインの状態変化の成功10/312件は，すべて参照列長0だった．非ゼロ長285件の成功は0件で，配置は0/103件だった．ゼロ長特例と実装条件の確認範囲の制約があるため，この集計を前提条件検証の独立した効果の根拠にはしない．設定の確認範囲は\ref{sec5.provenance}に示す．

前提条件検証の有無による成功率の差は，モデル・提示方式・タスク種別によって異なる．状態変化では検証ありの成功率が高いが，GPT-4o miniの配置タスクでは両提示方式で逆転する．

また，\ref{result1}と\ref{result2}では，配置タスクの成功率は状態変化タスクより全体的に低く，モデル間にも差がある．ただし，成功率の差だけから必要な推論能力や失敗過程を特定できない．対象の追跡，容器操作，終了判断などの具体的な要求は，後述する実行ログと代表ケースに分けて検討する．

\begin{table*}[tb]
\caption{全18条件の成功・失敗別行動列長（参照平均も各行と同じ集合で計算）}\label{tab:outcome_lengths}
\centering\small
\setlength{\tabcolsep}{3pt}
\begin{tabular}{llllrrrrrr}
\hline
タスク & 提示 & モデル & 検証 & 成功数 & 成功平均 & 参照平均 & 失敗数 & 失敗平均 & 参照平均 \\
\hline
状態 & Baseline & 4o & -- & 10 & 0.000 & 0.000 & 302 & 10.828 & 4.000 \\
状態 & Single & mini & あり & 195 & 3.969 & 3.662 & 117 & 6.043 & 4.222 \\
状態 & Single & 4o & あり & 246 & 3.923 & 3.927 & 66 & 6.091 & 3.667 \\
状態 & Multi & mini & あり & 234 & 4.179 & 3.735 & 78 & 7.128 & 4.282 \\
状態 & Multi & 4o & あり & 279 & 3.584 & 3.878 & 33 & 1.758 & 3.818 \\
状態 & Single & mini & なし & 192 & 4.203 & 3.792 & 120 & 5.942 & 4.000 \\
状態 & Single & 4o & なし & 237 & 3.979 & 3.932 & 75 & 5.933 & 3.680 \\
状態 & Multi & mini & なし & 217 & 4.558 & 4.028 & 95 & 6.274 & 3.516 \\
状態 & Multi & 4o & なし & 272 & 3.673 & 3.882 & 40 & 2.150 & 3.800 \\
配置 & Baseline & 4o & -- & 0 & -- & -- & 103 & 15.398 & 9.175 \\
配置 & Single & mini & あり & 48 & 8.625 & 8.771 & 55 & 6.109 & 9.527 \\
配置 & Single & 4o & あり & 76 & 8.868 & 9.329 & 27 & 5.852 & 8.741 \\
配置 & Multi & mini & あり & 42 & 7.929 & 8.548 & 61 & 5.426 & 9.607 \\
配置 & Multi & 4o & あり & 84 & 8.976 & 9.333 & 19 & 5.158 & 8.474 \\
配置 & Single & mini & なし & 49 & 7.490 & 8.286 & 54 & 5.778 & 9.981 \\
配置 & Single & 4o & なし & 70 & 8.671 & 9.029 & 33 & 4.455 & 9.485 \\
配置 & Multi & mini & なし & 45 & 7.578 & 8.400 & 58 & 5.017 & 9.776 \\
配置 & Multi & 4o & なし & 82 & 8.390 & 9.305 & 21 & 5.333 & 8.667 \\
\hline
\end{tabular}
\end{table*}

\begin{table*}[tb]
\caption{検証あり／なしの共通成功集合における平均行動列長}\label{tab:common_success_internal}
\centering\small
\setlength{\tabcolsep}{3pt}
\begin{tabular}{lllrrrr}
\hline
タスク & 提示 & モデル & 共通成功数 & 検証あり & 検証なし & 参照 \\
\hline
状態 & Single & mini & 161 & 3.826 & 3.919 & 3.627 \\
状態 & Single & 4o & 220 & 3.777 & 3.886 & 3.909 \\
状態 & Multi & mini & 186 & 4.027 & 4.220 & 3.796 \\
状態 & Multi & 4o & 267 & 3.584 & 3.678 & 3.888 \\
配置 & Single & mini & 42 & 8.119 & 7.571 & 8.333 \\
配置 & Single & 4o & 64 & 8.734 & 8.703 & 9.078 \\
配置 & Multi & mini & 36 & 7.694 & 7.306 & 8.000 \\
配置 & Multi & 4o & 78 & 8.872 & 8.397 & 9.346 \\
\hline
\end{tabular}
\end{table*}

\ref{tab:outcome_lengths}に成功・失敗別の記録列長を示す．失敗集合での短い列は途中終了を含む．参照平均は成功・失敗の各集合と同じタスクから計算し，件数0の平均は「--」とした．検証あり／なしの両方で成功したタスクにそろえた列長比較を表\ref{tab:common_success_internal}に示す．長さ比では参照列長0の状態変化27件・配置2件を分母から除き，件数を別記した．

\subsection{ログ付き追加比較（小規模診断）}\label{sec5.additional}

保存結果の再集計とは別に，2026年9月6--7日に，状態変化の\texttt{test\_task1, test\_task6}と配置の\texttt{test\_task1, test\_task65}を用いた追加比較を行った．4タスク，2モデル，10条件，各3反復の計240エピソードである．モデルは\texttt{gpt-4o-mini-2024-07-18}と\texttt{gpt-4o-2024-08-06}，temperatureは0，各要求の出力上限は256トークンとした．初期位置を固定し，各試行のAPI呼出し前に参照グラフとの記号状態一致，全ノードの位置差0.001以下，四元数距離0.0001以下を確認した．全240件がこの照合を通過した．これは過去の人物位置を復元した実験ではない．

C0は前提条件検証なし，C1は既存のLLM検証，C2は同じ抽出知識と条件辞書による記号検証である．C3は検証せず実行し，VH失敗後の最新知識とエラーに基づき1回修復する独自対照で，失敗実行も参照列長の2倍の実行予算に含める．CAPE等の忠実な再現とは位置付けない．C5は検証器だけに完全グラフを与え，計画器への知識は変えない．同じ8操作の条件を評価するが，既知の空状態集合と情報不足の区別もC2と異なるため，差を抽出漏れだけに帰属させない．

C1-budgetとC4-budgetは，毎回の行動決定で生成・レビュー・説明・再生成の4要求を行う専用対照である．前者は条件辞書，後者は条件辞書を提示しない一般的再考を用いる．両者の出力上限合計は1024トークンだが，入力長，実際の出力長，計画全体の要求数は一致しない．元の必要時のみ再生成するC1とは別仕様であり，厳密な総計算量一致は主張しない．R-random，R-fixed，R-noneはC1の事例だけを，保存済みの反復別ランダム1例，固定1例，空の事例メッセージに変更する．C1が意味類似検索の基準となる．反復ごとにタスク・モデルのブロックと条件順をseed 20260907で並べ替え，事前に保存した．

\begin{table}[tb]
\caption{追加比較の成功数（各条件4タスク$\times$3反復，分母12）}\label{tab:step4_pilot}
\centering\small
\begin{tabular}{lrr}
\hline
条件 & GPT-4o mini & GPT-4o \\
\hline
C0 & 8 & 9 \\
C1 & 9 & 9 \\
C2 & 4 & 10 \\
C3 & 5 & 9 \\
C1-budget & 4 & 11 \\
C4-budget & 6 & 10 \\
C5 & 5 & 9 \\
R-random & 8 & 9 \\
R-fixed & 7 & 9 \\
R-none & 3 & 3 \\
\hline
\end{tabular}
\end{table}

\ref{tab:step4_pilot}には参照再生が失敗した配置1件も含めた．別途，参照成功3件のみの集計も保存した．初期充足の状態変化1件は全条件・全反復で行動せず成功し，表の各分子に3件ずつ含まれる．4タスクは3種類の指示名に対応し，12個の独立タスクを評価したわけではない．C1とC0，C1-budgetとC4-budgetなどの差の方向はモデルによって異なり，一律な優位性や全415タスクへの一般化を示すものではない．

判定過程について，C1の候補に対するLLM判定を，抽出辞書上のJSON条件のtrue/falseと照合した正答数は，miniで34/43（誤受理4，誤棄却5），4oで45/54（誤受理9，誤棄却0）だった．これらは異なる軌跡で観測された判定であり，共通の固定問題集合に対するモデル比較ではない．情報不足のラベルは分母から除く規則としたが，このC1集計には該当例がなかった．

全10条件を合わせた終了判定はminiで496回，4oで600回であり，未達なのに厳密な文字列Endを返した回数はそれぞれ26回，0回だった．4oにはEnd/Continue以外の応答が2回あり，書式違反として別記した．一方，完全グラフ上のJSON条件がtrueだった実行でも，miniは376回中28回，4oは447回中37回でVHが失敗を返した．条件辞書の充足はVHの全実行制約の充足を保証しない．棄却して実行していない候補にVH成否を付与したものではない．

完全グラフ，抽出知識，全プロンプト・応答，生成候補，検証，再生成，実行前後の状態を保存した．新ドライバはEnd判定の入力となった直前グラフで成否を計算し，終了後にも独立に最終グラフを取得する．元NotebookのEnd後の取得時点と厳密に同一ではないが，今回の240件では成功ラベルと最終ゴール充足の不一致はなかった．成功API要求は3047件，入力1052576トークン，出力30594トークンであり，通常単価・キャッシュ割引なしの推定費用は1.681766米ドルだった．通信失敗と出力上限による打切りは0件である．設定・ソースのハッシュ，初期状態，ゴール，入力・応答，費用を照合した．個別値，反復別・対応比較，費用と再実行手順は補足資料\texttt{Paper/vh\_step4/}に示す．

\subsection{代表タスクのケース分析}\label{sec5.cases}
定量評価で見られる失敗とタスクに必要な判断を具体化するため，構成要素比較と同じ状態変化2件・配置2件を対象に，初期状態，ゴール，参照手順，実行過程を分析する．著者による説明的なケース分析である．以下はWindows VHの固定位置条件におけるケース分析用試行のログに基づき，構成要素比較の240件には合算しない．

状態変化\texttt{test\_task1}は，4個のlightswitch（71，173，261，427）が初期状態ですべてONであり，参照列長は0である．必要なのは指示の動詞に応じて操作を始めることではなく，対象全体の初期充足を確認して終了する判断である．

同じ指示の\texttt{test\_task6}では，173と427だけがOFFである．参照手順は173への移動とON操作，427への移動とON操作の4行動であり，ON済みの71と261との区別が必要となる．固定位置試行のGPT-4o mini・記号検証条件では，173への移動とON操作，261への移動が成功した後，427がOFFのまま終了判定がEndとなった．各行動の前提条件を満たすことと，タスク全体を達成することは異なる．この例では検証による再生成は起きていなかった．

配置\texttt{test\_task1}は，cupcake 195と196をkitchencounter 238に置く課題である．参照列は各物体への移動，把持，配置先への移動，配置を繰り返す8行動である．対象を2個とも追跡するとともに，保持・近接と物体の操作適合性を区別する必要がある．参照列の再生では最初の3行動が成功したが，4行動目のPUTが失敗した．保持・近接を確認できても実行可能性の十分条件とはならず，一般的なエラー応答だけから物理的原因を特定することはできなかった．またLLM検証条件では，OPEN deskという候補を棄却した後にOPEN cupcakeへ再生成し，VHが失敗を返す例があった．再生成後に再検証しない実装上の限界も確認できる．

配置\texttt{test\_task65}は，plum 53と54を初期CLOSEDのfridge 103に入れる課題である．参照列は1個目の移動・把持・冷蔵庫への移動・開扉・格納と，2個目の移動・把持・移動・格納の9行動であり，再生に成功した．途中の保持物と残り対象を管理する必要がある．検証なしのGPT-4o miniでは，53を格納し，54を右手に持ったままEndを返す例があった．保持情報は終了判定の入力にも含まれていたため，この失敗を情報抽出漏れだけでは説明できない．なお，両物体のINSIDE関係がゴールであり，最後に冷蔵庫を閉めることを評価時に追加要求してはいない．

これらは失敗過程と必要な判断を示す例である．4ケースの説明内容は2026年9月6日に確認を受けた．確認範囲は，本節の4ケースの初期条件・参照手順・必要な判断・難しい点の説明であり，確認者自身による実行や所要時間の計測は含まない．

行動列の実行可能性とゴール充足の関係を調べるため，基本性能評価の同じ4タスク・2モデルから記録された行動列8本を固定位置で再生した．8本とも記録列全体を再生でき，最終ゴールを満たしたのは5本だった．GPT-4o・Multiの配置\texttt{test\_task65}では10行動をすべて再生し，ゴールを満たした．これは\ref{sec5.replay}の棄却候補8件とは分析対象が異なる．記録されていない終端の失敗行動は含まず，生成時の人物位置との一致も確認していないため，記録接頭列を固定位置で実行した結果として解釈する．各列の長さとゴール充足は表\ref{tab:saved_replay_internal}に示す．


\subsection{考察}\label{sec5.5}
表\ref{tab:outcome_lengths}では成功・失敗ごとの手順長を，表\ref{tab:common_success_internal}では同じ成功タスク集合における検証有無の差を比較した．共通成功集合の状態変化では全4条件で検証ありの平均列長が短い一方，配置では全4条件で長かった．成功率に加えて，成功した手順の長さと途中終了を分けて評価する必要がある．各条件の分布，同一IDでの差・比，検索類似度別の詳細集計は公開リポジトリに保存する．
基本性能評価では配置の成功率が状態変化より全体的に低いが，タスク集合，参照列長，終了判定，VH実行制約も異なる．SingleとMultiの条件間では文面も異なる．

FRRMAは参照列長に基づく上限で終了したことを表すが，その直前にゴールへ到達した場合も含み，冗長な次行動を生成した証拠とは限らない．ETFRはEnd後のゴール不充足を表すが，抽出漏れ，候補選択，終了判断のどこに原因があるかは終端ラベルだけでは決まらない．ケース分析では保持情報が入力にあるのに未達終了した例があり，情報不足と入力の誤解釈を分ける必要性が示された．

構成要素比較のJSONラベルは8操作の条件辞書に対する診断である．完全グラフの条件がtrueでも実行失敗があり，LLMが棄却した候補には未実行のものが残る．SWITCHONのOFF条件は冗長操作を避ける条件にもなり得るが，条件充足・VH成否・目標達成・行動列長の改善はそれぞれ異なる指標である．

成功数に加え，修復機会と失敗実行数も解釈に関わる．C1とC3では途中の失敗実行数が異なる一方，C1-budgetとC4-budgetの比較は要求数と出力上限をそろえても総計算量はそろえていない．検索についても，対応比較では改善例と悪化例が共存する．これらの観測から，LLM逐次計画器の改善では，候補の条件判定，再生成の機会，未達ゴールを追跡する終了判断，条件辞書と実行制約の対応をそれぞれ検討する必要がある．保持情報が入力にある未達終了は終了判断の検討を，完全グラフ上の辞書条件がtrueでも生じる実行失敗は辞書と実行制約の対応の検討を促す．条件判定の一致率だけではこの二つを評価できない．本診断は，このタスク設定で観測した失敗を処理段階へ対応付け，改善対象を判断する材料を提供する．


\subsection{限界} \label{sec5.6}
第一に，データは2種類の定型指示と限られた物体クラスを扱い，事例・テスト間で家屋と操作構造を共有する．タスクの構築方針により，障害物移動や達成済みゴールを一時的に崩す後戻りを要する計画は評価対象外である．自由な言い換え，未知家屋，人間にとっての一般的な難易度も評価していない．構成要素比較は目的選択した4タスクの3反復であり，初期充足1件と参照再生失敗1件を含む．本比較の結果は，この4タスクの反復とVirtualHome上の実行条件の範囲で解釈する．

第二に，基本性能評価と構成要素比較では，タスク数，初期化の統制，実行記録の粒度が異なる．ベースラインは著者公開実装を変更せず，そのまま使用している．各評価の実行条件と記録範囲は\ref{sec5.provenance}に示す．C3は本実装上の独自対照である．計算量対照は要求数と出力上限をそろえた条件で行った．これらの制約から，一連の評価の各結果はそれぞれの比較範囲で解釈する．

第三に，グラフ取得後の抽出・自然言語化には情報の省略があり，条件辞書はVHの全制約を網羅しない．代表候補8件の接頭列方式で実行直前の状態照合を通過したのは2件で，人物のいないシーンへの復元方式で他の2件が通過したが，残る4件は本手順で状態を再現できず，状態照合を条件とする評価では候補の実行対象外だった．この4件も初回の再実行では実行済みで，2件がVH成功，2件がVH失敗だった．初回の実行成否と，後続評価の照合通過4件・未再現4件を区別して報告した．違反説明については，AI支援による基準説明・個別例照会後に著者の8件の記入を得た．要求文・応答・対応状態・採点理由を対応付けた小規模な事例評価として報告し，得られた著者記入結果の解釈は対象8件の範囲に限定する（\ref{sec5.review}）．

最後に，本研究はシミュレータ内部グラフと高レベル行動に依存し，視覚認識や実機制御を扱っていない．指定ゴール以外の原状回復，冷蔵庫の閉扉，容器の開放時間は追加の評価条件に含めない．
